\documentclass[conference]{IEEEtran}
\usepackage{cite}
\usepackage{amsmath,amssymb,amsfonts}
\usepackage{algorithmic}
\usepackage{graphicx}
\usepackage{textcomp}
\usepackage{xcolor}
\usepackage{booktabs}
\usepackage{multirow}
\usepackage{array}
\usepackage{url}
\def\BibTeX{{\rm B\kern-.05em{\sc i\kern-.025em b}\kern-.08em
    T\kern-.1667em\lower.7ex\hbox{E}\kern-.125emX}}

\begin{document}

\title{AdaCraft: Iterative Adaptive Personalization of Educational Examples via Agentic Feedback Loops}

\author{\IEEEauthorblockN{Anonymous Authors}}

\maketitle

\begin{abstract}
Most AI-powered educational systems treat personalization as a one-shot operation: a profile is read, an example is generated, and the system cannot adapt when learners signal dissatisfaction. We present \textbf{AdaCraft}, an agentic system that generates personalized educational examples through a multi-agent LangGraph workflow and refines them through natural-language feedback. The system operates through three capability layers: (1) a static user profile, (2) a Context Manager Agent that retrieves historical learning patterns and synthesizes a targeted personalization instruction before each generation, and (3) an Adaptive Response Agent that interprets free-form feedback and autonomously decides whether to regenerate, record a session insight, or persist a long-term learning pattern. We evaluate each layer's contribution via a four-tier ablation (T0: generic baseline, T1: +profile, T2: +context manager, T3: full system) grounded in LaMP \cite{b3} and scored on a five-axis rubric (Personalization Fidelity, Complexity Calibration, Conceptual Accuracy, Pedagogical Clarity, Domain Appropriateness) by a GPT-4.1-nano primary judge using G-Eval chain-of-thought \cite{b4} and Prometheus~2 rubric injection \cite{b5}, with inter-judge reliability validated against a Llama~3.3~70B secondary judge (Cohen's $\kappa \geq 0.607$). We further introduce three feedback-loop convergence metrics: FCR (Feedback Compliance Rate), LUR (Loop Utilization Rate), and PPU (Pattern Persistence Utilization), to measure iterative adaptation quality directly. Across two generator configurations (DeepSeek V3.2 and GPT-5-nano), the full system improves over the generic baseline by up to $+0.631$ for DeepSeek V3.2 and $+1.269$ for GPT-5-nano, with profile injection as the dominant driver in both runs. The agentic feedback loop achieves FCR@4 $\geq 0.891$, LUR $\geq 0.938$, and positive PPU deltas for returning users across both providers. Human evaluation with 10 participants (graduate students, postgraduate researchers, and practising educators) acting as learners in Study~A and as annotators in Study~B yields T3-preference rates of 90\% and 86.7\% respectively, validates LLM judge scores (composite human--LLM $r = 0.629$), and confirms multilingual robustness across six non-English languages (three providers, including Sarvam~105B on Indic languages) without generation quality degradation.
\end{abstract}

\begin{IEEEkeywords}
Natural language processing, Large language models, Multi-agent systems, Educational technology, Adaptive learning, Intelligent tutoring systems, Personalized learning, Artificial intelligence
\end{IEEEkeywords}


\section{Introduction}

Educational research has consistently shown that learners acquire and retain new knowledge most effectively when it is anchored in a familiar, concrete scenario, one that connects an unfamiliar abstraction to a domain the learner already inhabits \cite{b6}. The pedagogical value of an example, therefore, is inseparable from its relevance to the specific learner who encounters it.

This insight motivates the personalization problem in educational example generation. Large language models (LLMs) can now generate examples on demand for virtually any concept, but generating an example that is merely correct is not sufficient; it must be calibrated to the learner's background, domain familiarity, cultural context, and cognitive complexity preference. The gap between a correct example and a useful one is the gap that personalization must close.

Most LLMs and systems treat this as a one-shot problem: read the user profile once, generate a response, and stop \cite{b1, b2}. This is fundamentally at odds with how learning works. Learning is iterative and feedback-driven. Formative feedback, defined as information that helps a learner close the gap between where they are and where they need to be, is one of the most effective interventions in educational research \cite{b7}. Consider a nurse receiving an example that uses a pendulum to explain Newton's Second Law. She will likely signal, either explicitly or implicitly, that it missed the mark. A system that cannot act on this signal is not truly personalized; it is only initialized with a profile. The problem also extends beyond a single session: a learner's needs evolve over time, and a system with no cross-session memory cannot accumulate what it has observed about that user.

We address this gap with \textbf{AdaCraft}, an agentic system that generates contextually grounded examples on demand. The user interacts with it via a browser extension and provides natural-language feedback in their preferred language, and the system autonomously decides whether to regenerate the example with targeted modifications, record the interaction as a positive session insight, or persist a new behavioral pattern that will influence all future sessions for that user.

The core challenge is not generation quality in isolation but the design of an agentic architecture that can: (a) maintain a coherent personalization signal across sessions without manual curation, (b) interpret free-form feedback and act on it reliably, and (c) distinguish between transient session preferences and stable long-term traits worth persisting, without over-fitting to a single feedback signal at the expense of topic generalizability. We address them through a three-layer architecture and measure each layer's independent contribution via a four-tier ablation (Section~\ref{sec:evaluation}), where T0 is a generic baseline, T1 adds the user profile, T2 adds the Context Manager Agent, and T3 is the full system. The T1$\to$T2 delta isolates the value of cross-session context management, and the T2$\to$T3 delta isolates the contribution of feedback interpretation and long-term pattern persistence.

\textbf{Contributions.} This paper makes three contributions:
\begin{itemize}
    \item \textbf{AdaCraft}, an agentic system combining a static profile, a Context Manager Agent for cross-session pattern retrieval, and an Adaptive Response Agent for feedback-driven example refinement, deployable as a browser extension.
    \item \textbf{A four-tier ablation protocol} that isolates the independent contribution of each capability layer, grounded in the LaMP personalization benchmark and evaluated with a five-axis rubric across 256 cells and two generator configurations.
    \item \textbf{Three feedback-loop convergence metrics} (FCR, LUR, and PPU) that directly measure how well an iterative human-in-the-loop refinement system converges, complementing single-snapshot quality rubrics.
\end{itemize}

\section{Related Work}

The dominant paradigm in AI-powered personalized learning is static profile injection, where a user profile is prepended to the prompt \cite{b1,b2}. Benchmarks like LaMP \cite{b3} and LongLaMP \cite{b8} evaluate this approach but treat each generation as independent and open-loop, with no pathway for user dissatisfaction. While Self-Refine \cite{b12} explores iterative refinement using self-generated feedback, it lacks mechanisms for external human critiques and structured compliance metrics. Agentic frameworks such as LangGraph \cite{b9} provide workflow routing but leave feedback interpretation to the application layer. Furthermore, evaluation methods like G-Eval \cite{b4} and Prometheus 2 \cite{b5} assess single-shot quality but not feedback-loop convergence. No prior system jointly provides profile-grounded generation, feedback-driven in-session refinement with explicit agent decision logic, and cross-session pattern persistence. AdaCraft closes this gap and introduces feedback-loop convergence metrics (FCR, LUR, PPU) to measure iterative adaptation directly.

\section{System Architecture and Workflow}
\label{sec:system_workflow}

AdaCraft operates via a Chrome Extension interface and a Flask API backend hosting an agentic LangGraph workflow (Fig.~\ref{fig:workflow}). Upon installation, users define a static profile (role, background, learning style, complexity). When a user selects a concept to learn, the backend executes a multi-node workflow.

The workflow begins by injecting the user profile as a plain-English summary into the context. For returning users, a \textbf{Context Manager Agent} retrieves historical learning patterns and accepted insights, synthesizing a targeted personalization instruction. The \textbf{Generate Node} composes the example using the profile, context instruction, and any prior feedback, enforcing domain-appropriate structured outputs. 

After generation, the workflow pauses at the \textbf{User Review} node, displaying the example inline via the extension. If the user provides natural-language feedback, the \textbf{Adaptive Response Agent} activates. It autonomously selects from three tools: \textit{Regenerate Example} (issues a targeted rewrite directive and loops back to generation), \textit{Accept Example} (records a positive session insight), or \textit{Flag Pattern} (persists a long-term behavioral trait to influence future sessions). Cross-session data—including feedback history, learning patterns, and insights—are stored locally, enabling dynamic, long-term personalization without context-window overflow.

AdaCraft's dynamic personalization significantly improves over static baselines. For instance, given a generic baseline explaining compound interest via an agricultural metaphor, AdaCraft reframes the concept for a nurse in Lagos as a hospital equipment fund, using medical adherence analogies. When provided with a single natural-language critique (``This doesn't feel relevant to my field''), the Adaptive Response Agent autonomously reframes geological tectonic plate boundaries as medical anatomy analogies, demonstrating targeted, compliant refinement.

\section{Evaluation Design}
\label{sec:evaluation}

A central challenge in evaluating AdaCraft is that its value comes not just from the quality of a single example, but from how well the system improves across a feedback loop. Our evaluation design therefore has two parts: a four-tier ablation that isolates each capability layer's individual contribution, and a set of feedback-loop convergence metrics that characterize how the system evolves over multiple rounds.

\subsection{Ablation Structure}

We evaluate AdaCraft via a four-tier ablation grounded in the LaMP benchmark \cite{b3}, where each tier adds exactly one capability layer over the previous. This lets us answer three distinct questions: What does knowing the user's profile add? What does the Context Manager Agent add on top of the profile? What does the feedback loop and pattern memory add on top of both? Each tier is implemented via an evaluation-mode flag in the workflow state; four nodes perform conditional early-returns or skip disk writes based on this flag, requiring no graph rebuilds and ensuring the same codebase services all four conditions.

\begin{table}[htbp]
\caption{Ablation Tiers and Capability Layers}
\label{tab:tiers}
\begin{center}
\begin{tabular}{lccc}
\toprule
\textbf{Tier} & \textbf{Profile} & \textbf{Ctx.Mgr} & \textbf{Feedback \& Patterns} \\
\midrule
T0: Generic LLM   & --  & --  & --  \\
T1: T0+Profile      & \checkmark & --  & --  \\
T2: T1+ContextMgr   & \checkmark & \checkmark & -- \\
T3: Full AdaCraft  & \checkmark & \checkmark & \checkmark \\
\bottomrule
\end{tabular}
\end{center}
\end{table}

The T0$\to$T1 delta isolates the value of profile injection. T1$\to$T2 isolates the Context Manager Agent's contribution. T2$\to$T3 isolates the feedback loop and pattern memory.

\subsection{Evaluation Setup}

\textbf{Synthetic users.} We evaluate across 8 synthetic users, constructed as 4 cultural profiles $\times$ 2 start modes. The four profiles are: a high-school student (Berlin, European), a nurse (Lagos, West African), a humanities researcher (S\~{a}o Paulo, Latin American), and a software engineer (Tokyo, East Asian). This cross-cultural design tests whether the system grounds examples in culturally specific scenarios rather than falling back on generic templates.

\textbf{Topics.} Each user is tested on 4 topics drawn from distinct academic domains (Table~\ref{tab:topics}). The topics were chosen to stress different evaluation axes: psychology and earth sciences target cultural grounding and domain appropriateness; mathematics/finance targets complexity calibration; and evolutionary biology targets conceptual accuracy and pedagogical clarity.

\begin{table}[htbp]
\caption{Evaluation Topics and Primary Axes Tested}
\label{tab:topics}
\begin{center}
\begin{tabular}{lll}
\toprule
\textbf{Domain} & \textbf{Topic} & \textbf{Axis Tested} \\
\midrule
Evolutionary Biology & Natural Selection   & CA, PC \\
\textbf{Evaluation Scope.} We evaluate across 8 synthetic users (4 cross-cultural profiles $\times$ 2 start modes: cold vs. warm-start) and 4 academic topics, generating 128 cells per run (DeepSeek V3.2 and GPT-5-nano). Cold-start users test basic personalization, while warm-start users are pre-seeded with patterns to enable the PPU metric (T3 vs. T1 initial scores). For T3 sessions, we simulate three scripted feedback rounds: Round 1 delivers an \textit{easy} (unambiguous) or \textit{adversarial} (vague/contradictory) critique, Round 2 delivers acceptance, and Round 3 delivers a stable-trait statement to trigger pattern storage. T0--T2 run only one round, as they lack an active response agent.

\subsection{Five-Axis Evaluation Rubric}

Each example is scored 1--5 on five axes by an LLM judge using G-Eval chain-of-thought \cite{b4} (the judge reasons step by step before assigning a score) and Prometheus~2 rubric injection \cite{b5} (the full rubric definition is included in every judge prompt). Weights reflect a pedagogical priority ordering set before any data was collected, yielding the composite score in Eq.~\ref{eq:composite}.

\begin{equation}
C = 0.20 \cdot \text{PF} + 0.20 \cdot \text{CC} + 0.30 \cdot \text{CA} + 0.20 \cdot \text{PC} + 0.10 \cdot \text{DA}
\label{eq:composite}
\end{equation}

\begin{itemize}
    \item \textbf{Personalization Fidelity (PF, 0.20):} Does the example reflect the user's cultural background, occupation, and domain? Grounded in profile-conditioned generation benchmarks \cite{b3, b8}.
    \item \textbf{Complexity Calibration (CC, 0.20):} Does the example's depth and vocabulary match the user's requested complexity level? Motivated by learner-appropriate content design \cite{b13}.
    \item \textbf{Conceptual Accuracy (CA, 0.30):} Is the example factually and logically correct? Weighted highest because no degree of personalization can redeem a misconception.
    \item \textbf{Pedagogical Clarity (PC, 0.20):} Is the concept explained in a clear, structured, and learner-friendly way? Drawn from educational content quality frameworks \cite{b4, b13}.
    \item \textbf{Domain Appropriateness (DA, 0.10):} Does the example stay within the user's domain without inappropriate cross-domain elements? Weighted lowest as it largely follows from PF and CC in practice \cite{b13}.
\end{itemize}


\subsection{LLM Judge Design}

We use GPT-4.1-nano as the primary evaluation judge across both generator runs. \textbf{Run~1} uses DeepSeek V3.2 \cite{b14} as the generator; \textbf{Run~2} uses GPT-5-nano \cite{b15}. Keeping the judge fixed while varying the generator provides cross-provider evaluation independence and specifically mitigates self-enhancement bias \cite{b11}, where a model tends to score outputs resembling its own generation style more favorably, a concern noted for Run~2, where judge and generator share the same provider family. GPT-4.1-nano is cost-efficient at the scale of 256 evaluation cells (128 per run) while maintaining strong rubric-following capability. A secondary judge (Llama~3.3~70B) was deployed on a 20\% random subsample for inter-judge agreement verification (Cohen's~$\kappa$).

Two additional bias mitigations are applied:
\begin{itemize}

    \item \textit{Verbosity bias:} the judge system prompt explicitly instructs the model to penalize length and reward conciseness and clarity.
    \item \textit{Self-enhancement:} T0 generic examples are evaluated with separate judge invocations, not in direct comparison with profile-aware outputs, to prevent anchoring effects.
\end{itemize}

\subsection{Feedback Loop Convergence Metrics}

Standard quality rubrics measure a snapshot. For T3, we additionally compute three metrics that characterize the quality of the feedback loop itself over time. These metrics are designed to be applicable to any LLM system with a human-in-the-loop refinement process, not just AdaCraft.

\begin{itemize}
    \item \textbf{FCR (Feedback Compliance Rate):} The fraction of regenerations that correctly address the prior critique, as determined by the judge's compliance scoring. A regeneration is ``compliant'' if the judge assigns a compliance score $\geq 3$ on a 1--5 scale, the Likert midpoint, interpreted as ``the main substance of the critique was addressed even if imperfectly.'' We additionally report FCR@4 (score $\geq 4$) as a stricter quality threshold so that both an adequacy bar and a quality bar are visible and neither can be cherry-picked.
    \item \textbf{LUR (Loop Utilization Rate):} The fraction of T3 sessions where the Adaptive Response Agent triggered at least one regeneration in response to a complaint message. An LUR of 1.0 means the agent correctly identified every scripted complaint as requiring a revision, with no false negatives.
    \item \textbf{PPU (Pattern Persistence Utilization):} In short: do returning users with stored learning patterns get better first-generation examples than returning users without them? Formally, PPU is the mean PF score delta between warm-start T3 sessions (profile + Context Manager + seeded patterns) and warm-start T1 sessions (profile only, same users, same topics) on the \textit{initial} generated example, before any feedback is given. A positive delta indicates that the Context Manager Agent's retrieval and synthesis of stored patterns produces measurably better personalization than the static profile alone. Cold-start T3 sessions serve as a secondary control: if warm T3 PF $>$ cold T3 PF, the difference is attributable to the seeded patterns rather than to the Context Manager architecture itself.
\end{itemize}

\subsection{Statistical Testing Protocol}

To verify that the observed tier ordering reflects genuine improvements rather than sampling artefacts, we apply a two-stage non-parametric significance protocol to the composite scores of both generator runs. First, a \textbf{Friedman test} (a non-parametric repeated-measures analogue of one-way ANOVA) serves as an omnibus test across all four tiers. The Friedman test treats each (user, topic) pair as a block and ranks the four tier scores within it; a significant result rejects the null hypothesis that all four tiers are drawn from the same distribution. Second, all six pairwise tier combinations are tested with \textbf{Wilcoxon signed-rank tests} on the same paired cells, with \textbf{Holm-Bonferroni correction} applied across the six comparisons. We additionally report \textbf{rank-biserial $r$} as a sample-size-independent effect-size measure. Paired cells are defined by (user, topic), giving 32 pairs per run.

\section{Results and Discussion}
\label{sec:results}

We report results from two generator configurations evaluated on the same 128-cell grid, judged by GPT-4.1-nano. Run~1 uses DeepSeek V3.2 \cite{b14}; Run~2 uses GPT-5-nano \cite{b15}.

\subsection{Ablation Results}

Table~\ref{tab:ablation_results} reports mean composite and per-axis scores for each tier across both runs.

\begin{table}[htbp]
\caption{Ablation Results: Mean Scores by Tier (DS=DeepSeek V3.2, G5=GPT-5-nano)}
\label{tab:ablation_results}
\begin{center}
\scriptsize
\setlength{\tabcolsep}{6pt}
\begin{tabular}{llccccccc}
\toprule
\textbf{Run} & \textbf{Tier} & \textbf{PF} & \textbf{CC} & \textbf{CA} & \textbf{PC} & \textbf{DA} & \textbf{Comp.} & \textbf{$\Delta$T0} \\
\midrule
\multirow{4}{*}{DS}
 & T0 & 2.73 & 4.13 & 4.94 & 4.82 & 4.39 & 4.257 & {--} \\
 & T1 & 4.80 & 4.35 & 4.99 & 4.93 & 4.95 & 4.809 & +0.552 \\
 & T2 & 4.93 & 4.43 & 5.00 & 4.95 & 4.99 & 4.859 & +0.602 \\
 & T3 & 4.99 & 4.47 & 5.00 & 4.99 & 5.00 & 4.888 & +0.631 \\
\midrule
\multirow{4}{*}{G5}
 & T0 & 2.09 & 3.78 & 4.41 & 4.00 & 4.16 & 3.712 & {--} \\
 & T1 & 3.47 & 4.38 & 4.78 & 4.53 & 4.78 & 4.388 & +0.675 \\
 & T2 & 4.34 & 4.84 & 4.97 & 4.88 & 4.88 & 4.791 & +1.078 \\
 & T3 & 4.94 & 5.00 & 5.00 & 4.97 & 5.00 & 4.981 & +1.269 \\
\bottomrule
\end{tabular}
\end{center}
\end{table}

\begin{table}[htbp]
\caption{Pairwise Wilcoxon Signed-Rank Results (Holm-Bonferroni-corrected)}
\label{tab:significance}
\begin{center}
\scriptsize
\setlength{\tabcolsep}{4pt}
\begin{tabular}{lcccc}
\toprule
\textbf{Pair} & \multicolumn{2}{c}{\textbf{DeepSeek V3.2}} & \multicolumn{2}{c}{\textbf{GPT-5-nano}} \\
 & \textbf{$\Delta$} & \textbf{Rank-Biserial $r$} & \textbf{$\Delta$} & \textbf{Rank-Biserial $r$} \\
\midrule
T0$\to$T1 & $+0.552^{***}$ & 1.000 & $+0.675^{***}$ & 1.000 \\
T0$\to$T2 & $+0.602^{***}$ & 1.000 & $+1.078^{***}$ & 1.000 \\
T0$\to$T3 & $+0.631^{***}$ & 1.000 & $+1.269^{***}$ & 1.000 \\
T1$\to$T2 & $+0.049^{***}$ & 1.000 & $+0.403^{***}$ & 1.000 \\
T1$\to$T3 & $+0.078^{***}$ & 1.000 & $+0.594^{***}$ & 1.000 \\
T2$\to$T3 & $+0.029^{***}$ & 1.000 & $+0.191^{***}$ & 1.000 \\
\bottomrule
\multicolumn{5}{l}{$^{***}p < 0.001$ after Holm-Bonferroni correction. All 6 pairs significant in both runs.}
\end{tabular}
\end{center}
\end{table}

\subsection{Ablation Analysis}

Both generators produce a strict monotonic ordering T0 $<$ T1 $<$ T2 $<$ T3 (Table~\ref{tab:ablation_results}). The largest single gain is at T0$\to$T1 for both DeepSeek ($+0.552$) and GPT-5-nano ($+0.675$), confirming that profile injection dominates all other personalization signals. DeepSeek maintains a high baseline conceptual accuracy (T0 CA = 4.94) and requires less agentic layering to approach the ceiling (T0$\to$T3 total gain $+0.631$). Conversely, GPT-5-nano starts from a weaker baseline (T0 CA = 4.41, Comp = 3.712) but responds much more strongly to agentic layers, with a T1$\to$T2 gain ($+0.403$) substantially exceeding DeepSeek's, yielding the largest overall lift ($+1.269$). 

\textbf{Statistical significance.} The Friedman test rejects identical tier distributions ($\chi^2 \geq 91.35$, $p < 0.001$), confirming the T0~$<$~T1~$<$~T2~$<$~T3 ordering. All pairwise Wilcoxon signed-rank tests reach $p < 0.001$ after Holm-Bonferroni correction (Table~\ref{tab:significance}), with perfect monotonic improvement (rank-biserial $r = 1.000$). Inter-judge agreement with Llama 3.3 70B (20\% subsample) yielded $\kappa \geq 0.607$, confirming GPT-4.1-nano's reliability.

\subsection{Feedback Loop Convergence (T3)}

Table~\ref{tab:convergence} reports the three convergence metrics for T3 sessions.

\begin{table}[htbp]
\caption{T3 Feedback Loop Convergence Metrics (DS=DeepSeek V3.2, G5=GPT-5-nano)}
\label{tab:convergence}
\begin{center}
\scriptsize
\setlength{\tabcolsep}{6pt}
\begin{tabular}{llccc}
\toprule
\textbf{Metric} & \textbf{Description} & \textbf{DS} & \textbf{G5} & \textbf{Target} \\
\midrule
FCR@3 & Compliance $\geq 3/5$ & 1.000 & 0.945 & $\geq 0.75$ \\
FCR@4 & Compliance $\geq 4/5$ & 0.902 & 0.891 & $\geq 0.75$ \\
LUR   & Loop Utilization     & 0.938 & 0.969 & $\approx 1.0$ \\
PPU Delta PF & Warm T3 $-$ T1 PF & $+0.312$ & $+1.938$ & $> 0$ \\
\bottomrule
\end{tabular}
\end{center}
\end{table}

\textbf{FCR.} For DeepSeek V3.2: FCR@3 = 1.000 and FCR@4 = 0.902 (41 regenerations; mean compliance 4.34/5). For GPT-5-nano: FCR@3 = 0.945 and FCR@4 = 0.891 (55 regenerations; mean compliance 4.56/5). Both runs exceed the $\geq 0.75$ adequacy target by a wide margin. The easy/adversarial ordering differs across providers: DeepSeek scores higher on easy feedback (FCR@4 = 1.000) than adversarial (0.692), while GPT-5-nano inverts this (1.000 adversarial vs. 0.850 easy), suggesting that compliance difficulty is model-dependent. Overall FCR@4 exceeds the adequacy target in both runs, demonstrating that regenerations specifically address critiques rather than producing generic rewrites.

\textbf{LUR.} DeepSeek: LUR = 0.938 (30/32 sessions triggered at least one regeneration). GPT-5-nano: LUR = 0.969 (31/32). Both approach the $\approx 1.0$ target. On easy sessions (cold users, F1/F2), LUR = 1.000 in both runs --- the agent triggered regeneration on every unambiguous complaint. On adversarial sessions (warm users), LUR = 0.875 for DeepSeek and 0.938 for GPT-5-nano. The sub-unity adversarial LUR is expected in part by design: the near-zero-signal message (A4, ``Hmm.'') is intended to be classified as an acceptance rather than a regeneration request, so sessions receiving A4 correctly do not trigger regeneration. Remaining misses reflect genuinely ambiguous phrasing (A2: contradictory request) that the agent could not resolve into a regeneration decision.

\textbf{PPU.} DeepSeek: $\Delta$PF = $+0.312$ (warm T3 initial PF = 5.00 vs.\ warm T1 PF = 4.688; cold T3 PF = 4.979, $n = 16$ per group). GPT-5-nano: $\Delta$PF = $+1.938$ (warm T3 PF = 5.00 vs.\ warm T1 PF = 3.062; cold T3 PF = 4.875). In both runs, warm-start T3 sessions score higher on initial PF than warm-start T1 sessions (same users, same topics), confirming that stored learning patterns provide a meaningful personalization increment at first generation. For DeepSeek, the modest delta reflects that profile injection alone already achieves high personalization (PF $\approx$ 4.7--5.0), with patterns providing incremental gains. For GPT-5-nano, the larger delta (6.2x greater than DeepSeek) reflects its substantially lower baseline personalization (T1 PF = 3.062 vs DeepSeek T1 PF = 4.688), where patterns have significantly more room to improve initial examples, demonstrating differential value across generator types.


\section{Human Evaluation}
\label{sec:human_eval}

Two human studies complement the automated evaluation, with the same 10 participants (graduate students, postgraduate researchers, and practising educators from diverse disciplinary backgrounds) taking part in both. Study~A collects ecological feedback from participants evaluating examples on their own topics as learners; Study~B collects structured assessment on 30 examples drawn from the ablation corpus, with the same participants now acting as annotators.

\subsection{Study A: Learner Self-Evaluation}

Ten real learners from diverse backgrounds (engineering, medicine, social sciences, nursing, economics, education) each evaluated two T3 examples on their own topic of choice. Beyond the five rubric axes, participants scored Educational Usefulness (EU), Would Use (WU), Bias \& Fairness (BF), and expressed an Informed Preference (IP) between the T0 baseline and T3 final example. Table~\ref{tab:study_a} reports per-axis means. All three ecological indicators significantly exceed the neutral midpoint: IS~$t(19) = 4.945$, EU~$t(19) = 7.025$, WU~$t(19) = 4.723$ (all $p < 0.001$). T3 was preferred in 18/20 cases (90\%; binomial $p < 0.001$). The BF axis scored $4.80 \pm 0.41$ with all examples rated $\geq 4$, indicating no systematic bias detected by learners.

\begin{table}[htbp]
\caption{Study A: Learner Self-Evaluation ($n=20$, 10 learners)}
\label{tab:study_a}
\begin{center}
\scriptsize
\setlength{\tabcolsep}{3pt}
\begin{tabular}{lcccccccccc}
\toprule
 & \textbf{PF} & \textbf{CC} & \textbf{CA} & \textbf{PC} & \textbf{DA} & \textbf{FC} & \textbf{IS} & \textbf{EU} & \textbf{WU} & \textbf{BF} \\
\midrule
Mean & 4.05 & 3.85 & 4.45 & 4.20 & 4.00 & 3.95 & 4.15 & 3.95 & 3.90 & 4.80 \\
SD   & 0.69 & 0.59 & 0.51 & 0.62 & 0.73 & 0.83 & 1.04 & 0.61 & 0.85 & 0.41 \\
\bottomrule
\end{tabular}
\end{center}
\end{table}

\subsection{Study B: Contextual Evaluation by Annotators}

The same ten participants independently scored 30 ablation examples (16 cold-start, 14 warm-start; 15 DeepSeek, 15 GPT-5-nano) on eight axes: the five rubric axes plus FC, IS, and IP. Each annotator received the learner profile, T0 baseline, inter-round feedback history, and T3 final example, blind to provider. Table~\ref{tab:study_b} reports per-axis means, ICC(2,1) inter-rater reliability, and human--LLM Pearson correlations.

\textbf{IP \& IS.} T3 was preferred in 26/30 cases (86.7\%; binomial $p < 0.001$). IS mean $= 4.107$, significantly above neutral: $t(29) = 9.036$, $p < 0.001$.

\textbf{Inter-rater reliability.} ICC ranged 0.465--0.692, all $p < 0.001$. PF, CC, PC, FC, and IS reached moderate agreement (ICC $\geq 0.54$); CA and DA reached fair agreement, consistent with the difficulty of assessing correctness without deep domain expertise.

\textbf{Human--LLM agreement.} Composite $r = 0.629$ ($p < 0.001$); PF $r = 0.820$, CC $r = 0.728$, PC $r = 0.696$. DA showed near-zero correlation ($r = 0.025$, n.s.), reflecting dependence on domain expertise. FC human scores correlate strongly with automated FCR ($r = 0.748$, $p < 0.001$), validating FCR as a human-interpretable proxy. Warm-start examples score higher than cold-start on PF ($+0.226$) and IS ($+0.456$), providing human-judge confirmation of the PPU effect.

\begin{table}[htbp]
\caption{Study B: Annotator Evaluation — Mean Scores, ICC(2,1), Human--LLM $r$ ($n=30$, $k=10$)}
\label{tab:study_b}
\begin{center}
\scriptsize
\setlength{\tabcolsep}{4pt}
\begin{tabular}{lccccc}
\toprule
\textbf{Axis} & \textbf{Mean} & \textbf{SD} & \textbf{ICC} & \textbf{Interp.} & \textbf{$r_\text{h-llm}$} \\
\midrule
PF & 3.637 & 0.701 & 0.692 & Moderate & $0.820^{***}$ \\
CC & 3.867 & 0.582 & 0.596 & Moderate & $0.728^{***}$ \\
CA & 4.313 & 0.444 & 0.473 & Fair      & $-^\dagger$ \\
PC & 3.940 & 0.564 & 0.540 & Moderate & $0.696^{***}$ \\
DA & 4.197 & 0.395 & 0.465 & Fair      & 0.025 \\
FC & 3.707 & 0.578 & 0.566 & Moderate & $0.748^{***\ddagger}$ \\
IS & 4.107 & 0.671 & 0.684 & Moderate & --- \\
\bottomrule
\multicolumn{6}{l}{\scriptsize $^{***}p < 0.001$. $^\dagger$LLM CA near-ceiling. $^\ddagger$FC vs.\ FCR.} \\
\end{tabular}
\end{center}
\end{table}

\subsection{Multilingual Robustness}

The Adaptive Response Agent was tested in the FNL scenario without language-detection preprocessing across three provider configurations. DeepSeek V3.2 and GPT-5-nano were evaluated on six non-English languages; Sarvam~105B was evaluated on three Indic languages only (Hindi, Tamil, Bengali), as it does not support non-Indic scripts. DeepSeek achieves 0.882 mean routing accuracy (1.000 in Hindi and Bengali; 0.625 in German). GPT-5-nano and Sarvam both achieve uniform 0.667 across all tested languages, indicating model-dependent rather than language-dependent routing difficulty. Composite quality is maintained across all providers (DeepSeek: 4.725--4.900; GPT-5-nano: 4.700--4.900; Sarvam: 4.375--4.800), with no statistically significant degradation relative to English for any provider (Wilcoxon $p > 0.05$, Bonferroni-corrected $\alpha = 0.025$). Table~\ref{tab:multilingual_paper} summarises per-language accuracy.

\begin{table}[htbp]
\caption{FNL Routing Accuracy by Language (---\,=\,not evaluated)}
\label{tab:multilingual_paper}
\begin{center}
\scriptsize
\setlength{\tabcolsep}{3pt}
\begin{tabular}{lccc}
\toprule
\textbf{Language} & \textbf{DeepSeek V3.2} & \textbf{GPT-5-nano} & \textbf{Sarvam 105B} \\
\midrule
Hindi           & 1.000 & 0.667 & 0.667 \\
Tamil           & 0.833 & 0.667 & 0.667 \\
Bengali         & 1.000 & 0.667 & 0.667 \\
German          & 0.625 & 0.667 & ---   \\
Arabic          & 0.917 & 0.667 & ---   \\
Mandarin Chin.  & 0.917 & 0.667 & ---   \\
\midrule
\textbf{Overall} & \textbf{0.882} & \textbf{0.667} & \textbf{0.667}$^\dagger$ \\
\bottomrule
\multicolumn{4}{l}{\scriptsize $^\dagger$Indic languages only.}
\end{tabular}
\end{center}
\end{table}

\section{Limitations and Future Work}

\textbf{Synthetic user study.} The 8 evaluation users are constructed, not real learners. Although the profiles span four distinct cultural and professional backgrounds, actual users may reveal failure modes that scripted feedback cannot capture: for example, ambiguous feedback that the agent misinterprets, mid-session topic drift, or multi-issue feedback (e.g., ``make it simpler and use a different domain'') that the agent must decompose rather than act on as a single signal. A user study with real learners is planned as future work.

\textbf{Single-session pattern seeding.} The warm-start protocol seeds two patterns per user from a fixed template. Real pattern accumulation over many sessions may produce richer, more specific context instructions. A longitudinal study over at least 5 sessions per user would better characterize long-term PPU behavior.

\section{Conclusion}

AdaCraft shows that iterative, feedback-driven personalization is both architecturally feasible and measurably better than static profile injection. Across two generator configurations, DeepSeek V3.2 and GPT-5-nano, each capability layer adds a consistent, independent increment to personalization quality, with the full system (T3) outperforming a generic baseline by up to $+1.269$ composite points. The agentic feedback loop proves reliable in both runs: the Adaptive Response Agent correctly acts on user complaints at FCR@4 $\geq 0.891$ across both providers and regenerates in at least 93\% of sessions (LUR), and stored learning patterns measurably improve first-generation quality for returning users, with PPU deltas ranging from +0.312 (DeepSeek, already high-baseline) to +1.938 (GPT-5-nano, showing pattern impact at lower baseline). Human evaluation corroborates these findings: real learners prefer T3 in 90\% of cases (Study~A, $p < 0.001$) and the same participants acting as annotators prefer it in 86.7\% (Study~B, $p < 0.001$), composite human--LLM agreement reaches $r = 0.629$, and multilingual routing operates without degradation across six non-English languages. The four-tier ablation design and the FCR, LUR, and PPU metrics together offer a reusable evaluation framework for any LLM system that combines profile-grounded generation with iterative human-in-the-loop refinement, a paradigm that static personalization benchmarks are not designed to capture.

% Acknowledgment removed for blind review


\begin{thebibliography}{00}

\bibitem{b1}
Z. Zhang, R. A. Rossi, B. Kveton, Y. Shao, D. Yang, H. Zamani, F. Dernoncourt, J. Barrow, T. Yu, S. Kim, R. Zhang, J. Gu, T. Derr, H. Chen, J. Wu, X. Chen, Z. Wang, S. Mitra, N. Lipka, N. Ahmed, and Y. Wang, ``Personalization of Large Language Models: A Survey,'' \textit{arXiv preprint arXiv:2411.00027}, 2024.

\bibitem{b2}
J. Liu, Z. Qiu, Z. Li, Q. Dai, W. Yu, J. Zhu, M. Hu, M. Yang, T.-S. Chua, and I. King, ``A Survey of Personalized Large Language Models: Progress and Future Directions,'' \textit{arXiv preprint arXiv:2502.11528}, 2025.

\bibitem{b3}
A. Salemi, S. Mysore, M. Bendersky, and H. Zamani, ``LaMP: When Large Language Models Meet Personalization,'' in \textit{Proc. 62nd Annu. Meeting Assoc. Comput. Linguistics (ACL)}, Bangkok, Thailand, 2024, pp. 7370--7392.

\bibitem{b4}
Y. Liu, D. Iter, Y. Xu, S. Wang, R. Xu, and C. Zhu, ``G-Eval: NLG Evaluation Using GPT-4 with Better Human Alignment,'' in \textit{Proc. 2023 Conf. Empirical Methods Nat. Lang. Process. (EMNLP)}, 2023, pp. 2511--2522.

\bibitem{b5}
S. Kim, J. Suk, S. Longpre, B. Y. Lin, J. Shin, S. Welleck, G. Neubig, M. Lee, K. Lee, and M. Seo, ``Prometheus 2: An Open Source Language Model Specialized in Evaluating Other Language Models,'' in \textit{Proc. 2024 Conf. Empirical Methods Nat. Lang. Process. (EMNLP)}, 2024, pp. 4334--4353.

\bibitem{b6}
J. Sweller, ``Cognitive Load During Problem Solving: Effects on Learning,'' \textit{Cognitive Science}, vol. 12, no. 2, pp. 257--285, 1988.

\bibitem{b7}
J. Hattie and H. Timperley, ``The Power of Feedback,'' \textit{Rev. Educ. Res.}, vol. 77, no. 1, pp. 81--112, 2007.

\bibitem{b8}
I. Kumar, S. Viswanathan, S. Yerra, A. Salemi, R. A. Rossi, F. Dernoncourt, H. Deilamsalehy, X. Chen, R. Zhang, S. Agarwal, N. Lipka, C. V. Nguyen, T. H. Nguyen, and H. Zamani, ``LongLaMP: A Benchmark for Personalized Long-form Text Generation,'' \textit{arXiv preprint arXiv:2407.11016}, 2024.

\bibitem{b9}
LangChain~Inc., ``LangGraph: Build Resilient Language Agents as Graphs,'' 2024. [Online]. Available: \url{https://github.com/langchain-ai/langgraph}. [Accessed: Mar.~26, 2026].

\bibitem{b10}
B. Jiang, Z. Hao, Y.-M. Cho, B. Li, Y. Yuan, S. Chen, L. Ungar, C. J. Taylor, and D. Roth, ``Know Me, Respond to Me: Benchmarking LLMs for Dynamic User Profiling and Personalized Responses at Scale,'' in \textit{Proc. 2nd Conf. Lang. Modeling (COLM)}, 2025.

\bibitem{b11}
L. Zheng, W.-L. Chiang, Y. Sheng, S. Zhuang, Z. Wu, Y. Zhuang, Z. Lin, Z. Li, D. Li, E. P. Xing, H. Zhang, J. E. Gonzalez, and I. Stoica, ``Judging LLM-as-a-Judge with MT-Bench and Chatbot Arena,'' in \textit{Adv. Neural Inf. Process. Syst. (NeurIPS)}, vol. 36, 2023.

\bibitem{b12}
A. Madaan, N. Tandon, P. Gupta, S. Hallinan, L. Gao, S. Wiegreffe, U. Alon, N. Dziri, S. Prabhumoye, Y. Yang, S. Gupta, B. P. Majumder, K. Hermann, S. Welleck, A. Yazdanbakhsh, and P. Clark, ``Self-Refine: Iterative Refinement with Self-Feedback,'' in \textit{Adv. Neural Inf. Process. Syst. (NeurIPS)}, vol. 36, 2023.

\bibitem{b13}
K. K. Maurya, K. V. A. Srivatsa, K. Petukhova, and E. Kochmar, ``Unifying AI Tutor Evaluation: An Evaluation Taxonomy for Pedagogical Ability Assessment of LLM-Powered AI Tutors,'' in \textit{Proc. 2025 Conf. Nations Americas Chapter Assoc. Comput. Linguistics: Human Lang. Technol. (NAACL)}, 2025, pp. 1234--1251.

\bibitem{b14}
DeepSeek-AI, ``DeepSeek-V3 Technical Report,'' \textit{arXiv preprint arXiv:2412.19437}, 2024.

\bibitem{b15}
OpenAI, ``GPT-5-nano,'' OpenAI Platform Documentation, 2026. [Online]. Available: \url{https://platform.openai.com/docs/models}. [Accessed: Mar.~26, 2026].

\end{thebibliography}

\end{document}
